\documentclass[11pt,a4paper]{article}

\usepackage[utf8]{inputenc}

\usepackage[T1]{fontenc}

\usepackage[polish]{babel}

\usepackage{lmodern}

\usepackage{geometry}

\geometry{a4paper,margin=27mm}

\usepackage{microtype}

\usepackage[hidelinks]{hyperref}

% Portfolio: MSC-AR-P12; source: prospective ARC-XR research scope.

\title{Emocjonalne lustro XR wizualizujące oszacowany stan afektywny}

\author{}

\date{Wersja 0.1.0, 21 listopada 2026}

\begin{document}

\maketitle

\begin{abstract}

Przedmiotem opracowania jest Emocjonalne lustro XR wizualizujące oszacowany stan afektywny. Czy wizualizacja oszacowanego stanu emocjonalnego zmienia samoświadomość w porównaniu z neutralną prezentacją? W eksperymencie zostaną zestawione dwa sposoby informacji zwrotnej, wizualizacja estymowanego afektu oraz neutralny komunikat kontrolny. Głównym punktem końcowym jest zgodność samoopisu i późniejszego rozpoznawania własnych reakcji, bez utożsamiania z rozpoznaniem klinicznym. Opis ma charakter wstępny i nie przedstawia niezweryfikowanych wyników.

\end{abstract}

\section{Wprowadzenie}

Rozwój adaptacyjnych systemów XR stwarza możliwość badania interakcji człowieka z informacją przestrzenną, ale samo zastosowanie AR, VR lub modeli sztucznej inteligencji nie przesądza o przewadze poznawczej czy ergonomicznej. W rozpatrywanym zagadnieniu, Emocjonalne lustro XR wizualizujące oszacowany stan afektywny, kluczowe jest rozdzielenie właściwości systemu od zmienności zachowania użytkownika. Porównanie należy zaprojektować tak, aby mierzyć jednoznaczne zjawisko, a nie tylko ogólną ocenę technologii.

\section{Cel pracy i pytanie badawcze}

Czy wizualizacja oszacowanego stanu emocjonalnego zmienia samoświadomość w porównaniu z neutralną prezentacją? Projekt wymaga krytycznego przeglądu aktualnego stanu badań oraz zdefiniowania nowości w stosunku do innych tematów portfolio XR.

\section{Materiał i metody}

W eksperymencie zostaną zestawione dwa sposoby informacji zwrotnej, wizualizacja estymowanego afektu oraz neutralny komunikat kontrolny. Wymagane jest zdefiniowanie jednostki analizy, zasad doboru przypadków, wersji scenariusza i źródeł pomiaru. Rozkład warunków oraz reguły wykluczenia prób należy określić przed pomiarem. W badaniach z udziałem ludzi niezbędna jest uprzednia właściwa ocena etyczna i świadoma zgoda. W testach wyłącznie symulacyjnych należy oddzielić prawdę symulacyjną od obserwacji empirycznej.

\subsection{Wynik główny i kryteria jakości}

Zgodność samoopisu i późniejszego rozpoznawania własnych reakcji, bez utożsamiania z rozpoznaniem klinicznym. Pomiar powinien zachowywać jawne jednostki, znaczniki czasu i reguły oceny poprawności. Przy modelach predykcyjnych stosować rozłączne partycje oceny; przy eksperymentach interwencyjnych oddzielić skuteczność predykcji od przyczynowego wpływu interfejsu.

\section{Plan analizy}

Analiza wewnątrzosobnicza i jakościowa rozbieżności między oceną własną a sugestią interfejsu. W wynikach należy raportować wielkość efektu lub błąd estymacji z przedziałem niepewności, liczbą obserwacji i odsetkiem odrzuconych prób. Analizy dodatkowe wymagają jasnego oznaczenia jako eksploracyjne. Nie wybierać progów decyzyjnych ani modelu na podstawie obserwacji ze zbioru testowego.

\section{Wyniki}

W obecnej wersji nie raportuje się wyników pomiarów, tabel statystycznych ani parametrów wyznaczonych z nieudokumentowanych danych. To kierunek badawczy, a nie zatwierdzony protokół badania ani zgłoszony artykuł.

\section{Dyskusja}

Interpretacja przyszłych wyników musi uwzględnić alternatywne mechanizmy i zależność wniosków od jakości danych, warunków zadania oraz ograniczeń sprzętu. Wartość potencjalnego rozwiązania zależy nie tylko od trafności modelu, lecz także od jego opóźnienia, stabilności i wpływu na zachowanie człowieka. Wyniki testów syntetycznych nie mogą być automatycznie uogólniane na populację uczestników.

\section{Ograniczenia badania}

Interfejs może kierunkować samoopis; należy ujawnić niepewność predykcji i ograniczyć efekt sugestii. Granice nowości i trafności ekologicznej wymagają niezależnej weryfikacji. Żadna część szkicu nie dowodzi przeprowadzenia zbierania danych, uzyskania zgód etycznych ani gotowości manuskryptu do zgłoszenia.

\section{Wnioski}

Przedstawiono problem, kontrast i sposób jego planowanej oceny. Rozstrzygnięcie postawionego pytania wymaga przygotowania protokołu, zweryfikowanych źródeł oraz poprawnie udokumentowanych wyników. Bez tych materiałów nie można sformułować konkluzji empirycznych.

\end{document}

